% =====================================================================
%  Pb Action Portal -- Methodology & User Guide
%  The public, no-password Streamlit dashboard
%  Partnership for Battery Action (Pb Action / Global Development Incubator)
%
%  Compile with pdflatex (run twice for the table of contents / links):
%      pdflatex Portal_README.tex
%      pdflatex Portal_README.tex
% =====================================================================
\documentclass[11pt,a4paper]{article}

\usepackage[T1]{fontenc}
\usepackage[utf8]{inputenc}
\usepackage{lmodern}
\usepackage[margin=1in]{geometry}
\usepackage{xcolor}
\usepackage{graphicx}
\usepackage{booktabs}
\usepackage{longtable}
\usepackage{array}
\usepackage{enumitem}
\usepackage{titlesec}
\usepackage{fancyhdr}
\usepackage{amssymb}
\usepackage{amsmath}
\usepackage{listings}
\usepackage[hidelinks,colorlinks=true,linkcolor=pbblue,urlcolor=pbblue,citecolor=pbblue]{hyperref}

% ---- Palette -------------------------------------------------------------
\definecolor{pbblue}{HTML}{1E88E5}
\definecolor{pbgrey}{HTML}{5F6368}
\definecolor{pblight}{HTML}{F1F3F4}
\definecolor{pbgreen}{HTML}{43A047}
\definecolor{pborange}{HTML}{FB8C00}
\definecolor{pbrule}{HTML}{D0D3D6}

% ---- Section styling -----------------------------------------------------
\titleformat{\section}{\color{pbblue}\Large\bfseries}{\thesection}{0.6em}{}
\titleformat{\subsection}{\color{pbgrey}\large\bfseries}{\thesubsection}{0.6em}{}
\titleformat{\subsubsection}{\normalsize\bfseries}{\thesubsubsection}{0.6em}{}
\titlespacing*{\section}{0pt}{1.6em}{0.6em}

% ---- Header / footer -----------------------------------------------------
\pagestyle{fancy}
\fancyhf{}
\renewcommand{\headrulewidth}{0.4pt}
\renewcommand{\footrulewidth}{0pt}
\lhead{\small\color{pbgrey}Pb Action Portal}
\rhead{\small\color{pbgrey}Methodology \& User Guide}
\cfoot{\small\color{pbgrey}\thepage}

% ---- Table helpers -------------------------------------------------------
\newcolumntype{L}[1]{>{\raggedright\arraybackslash}p{#1}}
\renewcommand{\arraystretch}{1.25}

% ---- Code / value styling ------------------------------------------------
\lstset{
  basicstyle=\ttfamily\small,
  breaklines=true,
  columns=fullflexible,
  backgroundcolor=\color{pblight},
  frame=single,
  rulecolor=\color{pbrule},
  xleftmargin=0.5em,
  aboveskip=0.8em,belowskip=0.8em,
}
\newcommand{\code}[1]{\texttt{\small #1}}
\newcommand{\hs}[1]{\texttt{#1}}

% ---- Callout box ---------------------------------------------------------
\newcommand{\callout}[2]{%
  \vspace{0.4em}\noindent
  \fcolorbox{pbrule}{pblight}{%
    \begin{minipage}{0.97\linewidth}
    \small\textbf{\color{pborange}#1}\\[0.2em]#2
    \end{minipage}}\vspace{0.4em}}

% =====================================================================
\begin{document}

\begin{titlepage}
\centering
\vspace*{2cm}
{\color{pbblue}\Huge\bfseries Pb Action Portal\par}
\vspace{0.4cm}
{\Large The public lead-analysis dashboard\par}
\vspace{0.3cm}
{\large\color{pbgrey} Methodology \& User Guide\par}
\vspace{1.5cm}
{\large Partnership for Battery Action (Pb Action)\par}
{\normalsize Global Development Incubator (GDI)\par}
\vspace{2cm}
\begin{minipage}{0.8\textwidth}
\centering\small\color{pbgrey}
A public, interactive dashboard for exploring how lead moves through the global
economy --- international trade, mining and refining, battery manufacturing, and
end-of-life recycling. This document explains every tab, the calculations and
assumptions behind each view, and how to use and adjust the model's inputs. All
quantities are expressed in \textbf{metric tonnes of lead content} unless
otherwise stated.
\end{minipage}
\vfill
{\small\color{pbgrey} Document generated for the Pb Action team \textbullet\ contact
\href{mailto:ben.savonen@globaldevincubator.org}{ben.savonen@globaldevincubator.org}\par}
\end{titlepage}

\tableofcontents
\clearpage

% =====================================================================
\section{Overview}
% =====================================================================

The \textbf{Pb Action Portal} is the public Streamlit dashboard for the
Partnership for Battery Action. It brings together several third-party datasets ---
harmonized international trade statistics, mine and refinery production figures,
battery-collection data, and figures drawn from the published literature --- and
layers a set of interactive analyses on top of them, from bilateral trade maps to
a country-level material-flow model. It is built for high-level analysis and
pattern detection; every figure is expressed in tonnes of lead content so that
trade, production, and recycling flows can be compared on a common basis.

\subsection{No password; one-time disclaimer}

There is no login or password gate --- the app is public. The only interstitial is
a one-time data-quality disclaimer with an ``I understand --- enter the app''
button.

\callout{Data-quality disclaimer}{The Portal draws on large, multi-source
datasets --- BACI trade statistics, BGS/USGS production figures, and model-derived
estimates --- each with known gaps and reporting errors. Results are useful for
high-level analysis, rough estimates, and pattern detection, but should be
cross-checked against primary sources before drawing conclusions.}

\subsection{Flat tab structure}

The Portal presents \textbf{one flat level} of eight tabs:

\begin{longtable}{L{4.4cm} L{9.8cm}}
\toprule
\textbf{Tab} & \textbf{What it shows} \\
\midrule
\endhead
Trade Map            & Bilateral BACI lead-trade flows on a world map \\
Trade Trends         & Time series of imports/exports/production by country \\
Trade Relationships  & Flow-network graph of who trades with whom \\
Lead Accumulation    & Net lead balance (mining $+$ imports $-$ exports) over time \\
Production \& Capacity& Where lead is mined, refined, and where batteries are made \\
Recycling Economy Snapshot (Beta) & Country recycling-economy overview (radar) \\
Material Flow (Beta) & Experimental per-country material-flow Sankey \\
Literature Stats     & Searchable library of sourced datapoints \\
\bottomrule
\end{longtable}

\subsection{Units \& interpretation}

All values are in \textbf{metric tonnes of lead (Pb) content}; the Lead
Accumulation and model tabs display in kilotonnes (kt Pb). Model-derived numbers
are rounded to two significant figures for display.

\subsection{Data sources}

\begin{longtable}{L{5.5cm} L{8.5cm}}
\toprule
\textbf{Data} & \textbf{Source} \\
\midrule
\endhead
Bilateral trade flows & BACI / CEPII (harmonized UN Comtrade) \\
Mine \& refined production & USGS, BGS \\
Battery collection & Eurostat, BCI, ILZSG, national studies \\
Literature figures & Published studies (per-figure citations on the Literature
Stats tab) \\
\bottomrule
\end{longtable}

\subsection{Running \& deploying}

\begin{lstlisting}
pip install -r requirements.txt
streamlit run streamlit_app.py
\end{lstlisting}

No password is required. Dependencies are \code{streamlit}, \code{pandas},
\code{numpy}, \code{plotly}, and \code{pyvis} (for the interactive Flow Network).
All data the app needs ships in the repository, so no external downloads are
required. The Portal is deployed on Streamlit Community Cloud from the GitHub
repository \code{pb-action-portal}; it auto-redeploys on push to \code{main}.

% =====================================================================
\section{View mode: Easy vs.\ Advanced}
% =====================================================================

A \textbf{View mode} radio in the sidebar trades usability for control. It applies
to all data tabs (tabs 1--7).

\begin{itemize}[leftmargin=1.4em]
  \item \textbf{Easy mode} (default) fixes sensible assumptions and exposes only a
        \emph{Center year} slider. The fixed assumptions --- HS12 dataset, BGS
        source, 3-year centered average, default Pb factors with Slag and Other
        Lead Products excluded --- are restated in grey at the bottom of each tab.
  \item \textbf{Advanced mode} exposes every input: dataset (HS12/HS22), mining
        source (BGS/USGS), single-year vs.\ 3-year period, per-HS Pb-content
        factors, and extra per-tab controls (colour modes, layouts, animation,
        process-parameter sliders).
\end{itemize}

The Easy-mode factor set is constructed to exactly match Advanced mode at its
default widget states, so the two modes agree when Advanced is left at its
defaults.

\subsection{Datasets, sources, and time period}

\begin{itemize}[leftmargin=1.4em]
  \item \textbf{Dataset.} BACI from CEPII (harmonized UN Comtrade). \emph{HS12
        (2012--2024)} uses HS 2017 codes with waste-battery code \hs{854810};
        \emph{HS22 (2022--2024)} uses HS 2022 codes with the reclassified
        waste-battery code \hs{854911}.\footnote{\textbf{On using \texttt{854810}
        as a proxy for \texttt{854911}.} HS12's \texttt{854810} is a broad ``all
        spent batteries'' code, whereas HS22 introduced \texttt{854911}
        specifically for spent lead-acid batteries (ULABs). The historical
        (pre-2022) series therefore uses \texttt{854810} as a proxy for ULAB trade.
        Validated against the 2022--2024 overlap years (where BACI reports both
        nomenclatures), lead-acid makes up roughly 75--85\% of the spent-battery
        block by weight (77--85\% by value), confirming the proxy in direction and
        magnitude. However \texttt{854810} \emph{overstates} ULAB tonnage --- by
        about 10--25\% near 2022 --- because it also captures fast-growing lithium
        and other battery scrap; the proxy is cleaner (likely $>$90\%) the further
        back you go. Applying a scaling factor of roughly 0.90--0.95 for earlier
        years, tapering toward 0.80--0.85 near 2022, sharpens the historical series.
        (The validation is partly circular, since BACI's 2022--2024 \texttt{854810}
        is itself back-converted from HS22 reports.)} Easy mode is always HS12.
  \item \textbf{Mining \& refining source.} \emph{BGS} (1971--2023, total refined,
        no primary/secondary split) or \emph{USGS} (2015--2023, splits primary and
        secondary). When the selected source lacks a country/year, the app falls
        back to the other and shows a note.
  \item \textbf{Time period.} \emph{Single year}, or a \emph{3-year average}
        (recommended) over $[\,y-1,y,y+1\,]$ clipped to the data range; the
        average is annualised by summing flows and dividing by the number of years.
\end{itemize}

\subsection{Pb-content factors}\label{sec:pbfactors}

BACI reports gross product quantities; each HS code is multiplied by a Pb-content
fraction. In Advanced mode every factor is an adjustable slider (bounded lo--hi)
with a checkbox to exclude the code (factor $0.0$). The canonical table:

\begin{longtable}{l L{4.6cm} L{2.9cm} c c c}
\toprule
\textbf{HS code} & \textbf{Product} & \textbf{Category} & \textbf{Default} &
\textbf{Lo} & \textbf{Hi} \\
\midrule
\endhead
\hs{260700} & Lead Ore \& Concentrates      & Mining Outputs      & 0.600 & 0.45 & 0.75 \\
\hs{262021} & Leaded Gasoline Residues       & Slag                & 0.400 & 0.20 & 0.60 \\
\hs{262029} & Other Lead-Containing Residues & Slag                & 0.550 & 0.30 & 0.80 \\
\hs{282410} & Lead Monoxide                  & Battery Inputs      & 0.930 & 0.93 & 0.93 \\
\hs{282490} & Other Lead Oxides              & Battery Inputs      & 0.900 & 0.87 & 0.93 \\
\hs{780110} & Refined Unwrought Lead         & Battery Inputs      & 0.990 & 0.99 & 0.99 \\
\hs{780191} & Unwrought Lead (w/ Antimony)   & Battery Inputs      & 0.960 & 0.94 & 0.98 \\
\hs{780199} & Other Unwrought Lead           & Battery Inputs      & 0.970 & 0.95 & 0.99 \\
\hs{780200} & Lead Waste \& Scrap            & Battery Waste       & 0.725 & 0.50 & 0.95 \\
\hs{780411} & Lead Sheet/Strip/Foil (PCB)    & Other Lead Products & 0.990 & 0.99 & 0.99 \\
\hs{780419} & Lead Sheet/Strip/Foil (Other)  & Other Lead Products & 0.990 & 0.99 & 0.99 \\
\hs{780420} & Lead Tubes, Pipes \& Fittings  & Other Lead Products & 0.985 & 0.98 & 0.99 \\
\hs{780600} & Other Lead Articles            & Other Lead Products & 0.970 & 0.95 & 0.99 \\
\hs{850710} & Lead-Acid Starter Batteries    & New Batteries       & 0.600 & 0.55 & 0.65 \\
\hs{850720} & Other Lead-Acid Batteries      & New Batteries       & 0.600 & 0.55 & 0.65 \\
\hs{850790} & Battery Parts                  & Battery Inputs      & 0.625 & 0.30 & 0.95 \\
\hs{854810} & Used Lead-Acid Batteries (HS12)& Battery Waste       & 0.600 & 0.55 & 0.65 \\
\hs{854911} & Spent Lead-Acid Battery Waste (HS22) & Battery Waste & 0.600 & 0.55 & 0.65 \\
\bottomrule
\end{longtable}

\noindent\textbf{Default-off codes.} Slag (\hs{262021}, \hs{262029}) and Other
Lead Products (\hs{780411}, \hs{780419}, \hs{780420}, \hs{780600}) are excluded by
default; re-enable them in Advanced mode.

\callout{Where the sliders bite}{The trade tabs read a pre-computed
\code{actual\_lead} column (HS12's baked-in \code{Pb Quantity}, or a fixed
load-time factor $\times$ \code{Quantity} for HS22). The model tabs (Lead
Accumulation, Production BOTEC, Recycling Snapshot, Material Flow) re-derive lead
as \code{pb\_factors[code]} $\times$ \code{Quantity}, so the adjustable factors
visibly change those views.}

\subsection{Regions}

Places can be individual countries or regions. Region selectors list major
regions (Africa, Americas, Asia, Europe, Oceania) first, then UN sub-regions, then
countries A--Z. Regional trade views show cross-region (external) flows only.

% =====================================================================
\section{Trade tabs}
% =====================================================================

\subsection{Trade Map}

\textbf{What it shows.} A Plotly choropleth of BACI lead-trade for the selected
place, plus Top-30 partner tables. \emph{All countries} colours each country by
net position ($\text{exports}-\text{imports}$, signed-log normalised, red = net
importer, green = net exporter); a single country colours its partners; a region
colours its external partners (cross-region flows only).

\textbf{Calculations.} The single-country ``global volume scale'' multiplies a
balance ratio ($\text{net}/\text{total bilateral}$) by a log-intensity term so
small flows stay pale; the ``bilateral trade balance'' mode plots the ratio alone
(volume-independent), purple = pure receiver, orange = pure sender.

\textbf{Controls.} \emph{Select a country or region} (default All countries);
\emph{Color scale mode} (Advanced, non-All: Global volume scale / Bilateral trade
balance); \emph{Product selection} --- Easy shows the five trade-category
checkboxes, Advanced per-HS checkboxes grouped by six material-flow categories
(Slag and Other Lead Products off by default).

\subsection{Trade Trends}

\textbf{What it shows.} Line charts of import/export volume (or production) over
time, for one place (Easy) or up to three (Advanced).

\textbf{Calculations.} Trade series sum \code{actual\_lead} by year for the chosen
category's HS codes (raw annual values --- not subject to the 3-year window).
Production series come from BGS or USGS with automatic source fallback. The
Relative view plots year-on-year \% change.

\textbf{Controls.} \emph{Country / Region} slots; \emph{Series} multiselect
(default Total Lead-Related Imports \& Exports); \emph{View} absolute/relative
(Advanced); \emph{Shared y-axis} toggle (Advanced).

\subsection{Trade Relationships (Flow Network)}

\textbf{What it shows.} A node-link graph: bubble size
$\propto\sqrt{\text{total trade volume}}$, arrow width
$\propto\sqrt{\text{flow volume}}$, arrows coloured by category. An interactive
draggable view (default) or a static Plotly figure.

\textbf{Controls.} \emph{View mode} (Advanced); \emph{Select countries} (default
India, China, USA); \emph{Minimum flow} (Advanced 100--10000, default 500; Easy
fixed 100); \emph{Layout} (default geographic centroids; Advanced adds Grid and
Force-directed); \emph{Focal country}; \emph{Region cluster}, \emph{Flow
direction}, \emph{Use all years}, and \emph{Hide country pairs} (all Advanced);
product selection as on the Trade Map.

% =====================================================================
\section{Model \& material-flow tabs}
% =====================================================================

\subsection{Lead Accumulation}

\textbf{What it shows.} A cumulative net-lead line chart since the 2012 base year,
a per-category net bar for the selected year, and a diverging world map of every
country's net balance (Advanced adds an annual bar, metric tiles, and an HS-code
table).

\textbf{Calculation.}
\[
\text{Net} \;=\; \text{Mining} \;+\; \text{Imports}_{\text{Pb}}
\;-\; \text{Exports}_{\text{Pb}},
\]
where mining is mine production plus net ore trade (\hs{260700}), and trade is
converted with the adjustable Pb factors (unknown codes default $0.70$). Positive
= net accumulator, negative = net depleter. The world map uses a diverging
red--green scale centred at zero, clipped at the 95th percentile of magnitude.

\textbf{Controls.} \emph{Select country or region} (default Ghana); inherits the
sidebar year, dataset, Pb factors, and mining source.

\callout{A trade-based approximation}{Net balance does not account for domestic
consumption, stockpiles, or informal flows; 2024 data is incomplete.}

\subsection{Production \& Capacity}

\textbf{What it shows.} A log-scale choropleth of one of four datasets, with a
ranked table for reported datasets and (Advanced) an animated time-lapse.

\begin{longtable}{L{5.4cm} L{2.0cm} L{6.1cm}}
\toprule
\textbf{Dataset} & \textbf{Type} & \textbf{Source / metric} \\
\midrule
\endhead
Lead Mining                    & reported & BGS or USGS mine production \\
Lead Refining                  & reported & BGS total refined, or USGS primary $+$
secondary \\
Battery Manufacturing (BOTEC)  & model    & \code{F4\_batt\_lead} (see engine) \\
Battery Consumption (BOTEC)    & model    & \code{F5\_implied} (see engine) \\
\bottomrule
\end{longtable}

\textbf{Calculations.} Reported datasets are pulled from an outer-join of separate
USGS and BGS extracts (never reconciled). The two BOTEC datasets run the
material-flow engine (Section~\ref{sec:engine}) once per eligible country.

\textbf{Controls.} \emph{Dataset} radio; \emph{Animated} toggle (Advanced);
\emph{Year} slider (Advanced); a page-local BGS/USGS radio in Easy mode; a
\emph{Process parameters} expander for the BOTEC datasets in Advanced.

\subsection{Recycling Economy Snapshot (Beta)}

\textbf{What it shows.} Side-by-side radar (``diamond'') charts comparing
countries (up to three in Advanced, two in Easy) on four recycling axes ---
Collection, Breaking, Secondary Smelting, Manufacturing (optionally plus Mining).
A perfectly circular economy plots as an equilateral diamond.

\textbf{Calculation.} Axis values are material-flow engine outputs
(\code{B4\_collected}, \code{B2\_break\_out}, \code{eff\_secondary},
\code{F4\_batt\_lead}, \code{mined}) computed with \emph{default} process
parameters --- the Advanced sliders do not apply here. An optional Collection
Baseline marks the ``perfect loop'' reference.

\textbf{Controls.} Country selectboxes (defaults Germany, USA); Advanced toggles
for Mining axis, Collection Baseline, and Same scale. In Easy mode the baseline is
always on.

\subsection{Material Flow (Beta)}

\textbf{What it shows.} A per-country material-flow Sankey (collection $\to$
breaking $\to$ secondary smelting $\to$ manufacturing $\to$ installation), with
per-commodity trade nodes and (Advanced) a model-output table.

\textbf{Controls.} \emph{Select country} (restricted to countries with refining
data); \emph{Min flow threshold} (Advanced, default 10 t Pb; Easy fixed 10); the
full \emph{Process parameters} expander in Advanced. Advanced draws gross
import/export nodes; Easy nets each commodity unless both directions exceed the
threshold.

\subsection{The material-flow engine}\label{sec:engine}

Four of the tabs are powered by one underlying model: the \textbf{Material Flow}
Sankey, the \textbf{Recycling Economy Snapshot}, and the two BOTEC datasets on
\textbf{Production \& Capacity} (Battery Manufacturing and Battery Consumption).
Understanding this engine explains most of what those tabs show.

\paragraph{What it is --- in plain terms.} For a single country and year, the model
reconstructs the whole lead cycle --- from batteries being collected at the end of
their life, through recycling, all the way to new batteries being made and put into
service. It works by starting from one quantity we can actually measure (how much
lead that country's smelters produced from scrap) and reasoning outward in both
directions: trade data accounts for what crosses borders, and a short list of
efficiency assumptions accounts for the losses at each step. It is a single,
transparent pass of arithmetic --- not a statistical fit --- so every output can be
traced straight back to its inputs.

\paragraph{The anchor.} The one firmly-measured quantity is a country's
\textbf{secondary (scrap-based) smelting output}, reported by BGS or USGS.
Everything else is estimated relative to it. Anchoring on a real, reported number
keeps the whole chain tied to observed production rather than to a tower of
assumptions. (A country that reports no refining at all cannot be anchored and is
skipped.)

Two pieces of notation below: for each traded product, ``imports'' and ``exports''
are that country's trade in it (in tonnes of lead content), and the Greek letters
are the efficiency/share parameters listed in the table further down.

\paragraph{Working backward --- from the smelter up to collection.} Moving
\emph{up} the recycling chain, each step asks: how much had to enter this stage for
the next stage's figure to come out? That means dividing by an efficiency (to undo
the losses) and adjusting for net trade of the material at that stage.
\begin{align*}
\text{anchor} &= \text{reported secondary smelting output} \\
B_1\ (\text{scrap into smelters}) &= \text{anchor} / \eta_{\text{secondary}} \\
B_2\ (\text{lead recovered by breaking}) &= \max(0,\; B_1 - \text{imports}_{\text{scrap}} + \text{exports}_{\text{scrap}}) \\
B_3\ (\text{used batteries broken}) &= B_2 / (\delta \cdot \eta_{\text{break}}) \\
B_4\ (\text{used batteries collected}) &= \max(0,\; B_3 - \text{imports}_{\text{used}} + \text{exports}_{\text{used}}) \\
B_5\ (\text{end-of-life batteries generated}) &= B_4 / \gamma \\
B_6\ (\text{uncollected / disposed}) &= B_5 \cdot (1-\gamma)
\end{align*}
In words: the scrap the smelters consumed ($B_1$) must have come from domestic
battery-breaking plus net imported scrap ($B_2$); that recovered lead implies a
larger quantity of used batteries fed in, before breaking and end-of-life losses
($B_3$); those came from batteries collected at home plus net trade in used
batteries ($B_4$); and the collected amount is only a fraction $\gamma$ (the
collection rate) of all the batteries that actually reached end of life ($B_5$),
the remainder being lost or dumped ($B_6$).

\paragraph{Working forward --- from feedstock down to new batteries.} From the
total refined metal available, the model moves \emph{down} the manufacturing chain:
\begin{align*}
\text{feedstock} &= \text{primary metal} + \text{secondary metal} + \text{net refined \& crude imports} \\
F_2\ (\text{lead into non-battery uses}) &= \text{feedstock} \cdot (1-\beta) \\
F_3\ (\text{lead routed to batteries}) &= \text{feedstock} \cdot \beta \\
F_4\ (\text{lead in batteries made}) &= \max(0,\;F_3)\cdot\eta_{\text{mfg}} \\
F_5\ (\text{implied battery consumption}) &= F_4 + \text{imports}_{\text{batt}} - \text{exports}_{\text{batt}}
\end{align*}
In words: the metal a country has to work with (feedstock) is its own refined
output plus net imports of refined feedstock and crude lead (crude is first
de-rated by a refining recovery $\eta_{\text{refine}}$); a share $\beta$ of that
goes into batteries and the rest into other lead products ($F_2$/$F_3$);
manufacturing losses reduce it to the lead actually embodied in new batteries
($F_4$); and adding net battery trade gives the lead in batteries \emph{put into
service} domestically ($F_5$).

\paragraph{Where the primary/secondary split comes from.} The anchor has to be
divided between metal made from ore (primary) and metal made from scrap
(secondary). Estimated primary metal is domestic ore
($\text{mined}+\text{imports}_{\text{ore}}-\text{exports}_{\text{ore}}$) times a
smelting recovery $\eta_{\text{ore}}$; if that comes to less than 5\% of reported
refined output, the model treats \emph{all} refined lead as secondary. When USGS
data is selected, its reported primary/secondary split is used directly instead.

\callout{Key assumptions --- read before trusting a number}{The engine is a
back-of-the-envelope (BOTEC) estimate, deliberately simple. In particular it
assumes: (1)~\textbf{steady state} --- a single trade year is treated as if the
system were in balance, with no time lag between batteries sold and batteries
returned; (2)~\textbf{one aggregate flow} --- there is no split between formal and
informal recycling, which in reality can change the efficiencies at every step;
(3)~\textbf{global default efficiencies} --- the recovery and collection rates are
literature midpoints applied to every country until you change them;
(4)~\textbf{trade as a clean proxy} --- traded scrap, used batteries, and feedstock
are assumed to be correctly classified and measured, so any gaps or
misclassifications in the underlying trade data flow straight through; and
(5)~\textbf{no reconciliation} --- it is a single deterministic pass, not a model
solved for overall consistency against independent stock data. Outputs are rounded
to two significant figures and labelled Beta in-app.}

\paragraph{Adjustable parameters (Advanced mode).} Each stage's efficiency or share
is a slider. The defaults are neutral starting points, \emph{not} country-specific
values.

\begin{longtable}{L{6.4cm} c c c}
\toprule
\textbf{Parameter} & \textbf{Symbol} & \textbf{Range} & \textbf{Default} \\
\midrule
\endhead
Secondary smelting recovery & $\eta_{\text{secondary}}$ & 0.80--1.00 & 0.97 \\
Breaking recovery           & $\eta_{\text{break}}$     & 0.70--1.00 & 0.95 \\
Pb retained at end-of-life  & $\delta$                  & 0.80--1.00 & 0.95 \\
Battery share of lead demand& $\beta$                   & 0.50--1.00 & 0.85 \\
Manufacturing efficiency    & $\eta_{\text{mfg}}$       & 0.90--1.00 & 0.98 \\
Primary smelting recovery   & $\eta_{\text{ore}}$       & 0.80--1.00 & 0.95 \\
Refining recovery (crude lead) & $\eta_{\text{refine}}$ & 0.80--1.00 & 0.97 \\
Collection rate             & $\gamma$                  & 0.30--1.00 & 0.95 \\
\bottomrule
\end{longtable}

\callout{Start from the defaults, then tune for each country}{Every parameter ships
with a single \emph{global} default (for example, the collection rate
$\gamma=0.95$), applied identically to every country. Real-world values differ
enormously: collection rates alone range from roughly 0.99 in the USA and Japan, to
about 0.70 in India, to around 0.60 in Nigeria --- and the battery share of lead
demand, breaking and smelting recoveries, and end-of-life retention vary just as
much. The out-of-the-box numbers are therefore \emph{illustrative}, not a
country-specific answer. For any serious analysis, switch to Advanced mode, open
\emph{Process parameters}, and set at least the collection rate ($\gamma$), the
battery share ($\beta$), and the recovery efficiencies to values appropriate for
the country and year you are studying. The Sankey and snapshot update immediately,
so you can see how sensitive the picture is to each assumption.}

% =====================================================================
\section{Literature Stats}
% =====================================================================

\textbf{What it shows.} A searchable, read-only library of quantitative datapoints
from the lead-battery literature, each tied to a source citation (reading
\code{data/literature\_datapoints.csv}). Internal model-variable codes are hidden
and mapped to plain-English topics (Collection \& ULAB volumes, Recycling
efficiency, Battery composition, Lead demand, Trade flows, Health \&
environmental exposure, and so on).

\textbf{Controls.} A \emph{Search} box; \emph{Topic}, \emph{Geography}, and
\emph{Evidence type} multiselects; a \emph{Year} slider with an ``Include undated /
multi-year'' toggle. Results appear in a sortable, single-select table; clicking a
row opens a detail panel with the verbatim quote and APA citation. A CSV download
exports friendly columns. The only aggregation shown is the count of matching
statistics.

\callout{Submission form is built but disabled}{The Portal contains a full
visitor-submission form (\code{literature/submit.py}) that would write new
datapoints to a Google Sheet for review, with an offline CSV/mailto fallback. It
is currently \emph{disabled}: the \code{render\_submission\_form()} call and its
import are commented out in \code{literature/app.py}, and \code{gspread} /
\code{google-auth} are commented out in \code{requirements.txt}. To re-enable it:
uncomment those lines, restore the two dependencies, and configure the secrets
(see \code{.streamlit/secrets.toml.example}).}

\end{document}
